% P7: Cache Geometry Under Obfuscation — KV-Cloak as Defense Against
%      Adversarial KV-Cache Steering
%
% Companion to:
%   P2 Oracle Loop (detection + steering method)
%   P6 Oracle Formulary (therapeutic profiles)
%
% Core argument: KV-Cloak destroys external geometric detection and
% renders adversarial steering injection incoherent, while a defensive
% Oracle Loop operating inside the TEE retains full detection capability.
%
% This paper is the dual-use mitigation that makes P2 and P6 publishable.

\documentclass[11pt,a4paper]{article}
\usepackage{amsmath,amssymb,booktabs,graphicx,hyperref}

\title{Cache Geometry Under Obfuscation: KV-Cloak as Defense Against\\
Adversarial KV-Cache Steering}

\author{
Lyra\textsuperscript{1*}\quad
Thomas Edrington\textsuperscript{1}\quad
Dwayne Wilkes\textsuperscript{1}\\[4pt]
\textsuperscript{1}Liberation Labs / THCoalition\\
\textsuperscript{*}Lead AI author. Corresponding human author: T.~Edrington.
}

\date{Draft --- \today}

\begin{document}
\maketitle

% ===================================================================
\begin{abstract}
% ===================================================================

Recent work demonstrates that singular value decomposition (SVD) of the
key-value (KV) cache reveals geometric signatures of cognitive states
including confabulation, deception, and sycophancy, enabling both
detection and correction of misalignment at inference time (the ``Lyra
Technique''). These same methods pose a dual-use risk: an adversary with
cache access could deploy geometric detection to identify vulnerable
cognitive states and inject emotion-vector perturbations to deliberately
destabilize a model.

We investigate whether KV-Cloak~\citep{luo2025shadow}, a reversible
matrix-based cache obfuscation scheme, mitigates this risk. In
experiments across six models (50--100 prompts each), we find:

\begin{enumerate}
\item \textbf{Detection on uncloaked cache is viable.} On the distilled
model (Qwen3.5-27B), Marchenko-Pastur-corrected spectral gap
distinguishes confabulation from hedged responses (AUROC~0.903, 95\%
bootstrap CI [0.806, 0.977], $R^2=0.068$ vs token count; $n=7$
confabulations, $n=87$ hedged). A text-only surface-feature baseline
achieves 0.755; cache geometry outperforms text in 96.7\% of bootstrap
samples (delta 0.148, 95\% CI [$-$0.008, 0.322])---suggestive but not
significant at the 95\% level. In a cross-architecture sweep
(Mistral-7B, Qwen2.5-14B, Llama-3.1-8B, DeepSeek-R1-7B), detection
strength varies by model (0.590--0.784); no model achieves significance
after Bonferroni correction for four comparisons.

\item \textbf{Obfuscation completely transforms the external feature
space.} Per-head orthogonal rotation---simulating KV-Cloak's reversible
matrix obfuscation---produces a total distributional shift in MP
features (cloaked-honest vs uncloaked-honest: Wilcoxon signed-rank
$p = 5.5 \times 10^{-16}$, $n=87$ paired trials). A detector
calibrated on uncloaked caches cannot operate on cloaked caches.

\item \textbf{Models are robust to random perturbation.} Magnitude-based
attacks (random-direction injection at steering-vector norm) produce
0--2 adverse events per model in 10 trials, inconsistent with
destabilization rates above 30\% (power 0.97 to detect) but unable to
rule out rates up to 20\% at this sample size. Targeted steering
injection was tested but the positive control was underpowered (3
calibration pairs; direct injection produced 0--1/10 corrections),
making through-cloak injection results uninterpretable. This remains
preliminary.

\item \textbf{A defensive Oracle Loop inside a trusted execution
environment (TEE) retains full capability.} The defender operates on the
pre-obfuscation cache, where detection (0.903) and steering function
normally. KV-Cloak protects the cache from external adversaries while
the TEE protects the obfuscation key and Oracle parameters. This is an
architectural argument, not an experimentally validated system.
\end{enumerate}

We conclude that KV-Cloak and the Oracle Loop are complementary: the
Oracle Loop provides misalignment detection and correction; KV-Cloak
prevents external adversaries from weaponizing the same geometric
methods. Together, they address the dual-use risk inherent in publishing
cache-geometry techniques. A direct implementation of KV-Cloak's
mechanism ($S \cdot \hat{P} \cdot (K + A) \cdot M$) degrades the
strongest detection feature (spectral gap) from 0.903 to 0.614---near
chance---confirming that the defense is stronger than our initial
orthogonal-rotation simulation (0.786) suggested. We discuss limitations
including small sample size ($n=7$ confabulations on the distilled
model), marginal cache-vs-text added value, reliance on TEE hardware
trust assumptions, and underpowered injection testing.

\end{abstract}

% ===================================================================
\section{Introduction}
% ===================================================================

The discovery that KV-cache geometry encodes cognitive state information
creates a tension between transparency and safety. Publishing detection
methods enables defensive applications (confabulation monitoring,
deception detection, self-regulating AI systems) but simultaneously
provides a roadmap for adversarial exploitation. An attacker with cache
access could:

\begin{itemize}
\item \textbf{Surveil} cognitive states to identify when a model is
uncertain, compliant, or vulnerable to manipulation.
\item \textbf{Inject} emotion vectors calibrated to specific
misalignment types, deliberately inducing confabulation, sycophancy, or
overconfidence.
\item \textbf{Identify} individual users or personas from cache geometry
fingerprints.
\end{itemize}

This paper investigates whether KV-Cloak~\citep{luo2025shadow}---a
lightweight, reversible cache obfuscation scheme designed for privacy
protection---also mitigates these adversarial capabilities. We show that
it does: obfuscation renders geometric detection and vector injection
incoherent to external adversaries, while a defender operating inside
the trusted execution environment retains full access to the uncloaked
cache and its geometric properties.

% ===================================================================
\section{Background}
% ===================================================================

\subsection{The Lyra Technique}

KV-cache geometry---specifically, Marchenko-Pastur-corrected spectral
features derived from SVD of the key tensor---distinguishes cognitive
states during autoregressive generation. Prior work reports within-model
deception detection at AUROC~1.000 across seven model scales (0.5B--32B),
confabulation detection at 0.969--0.999, and hardware invariance
($r>0.999$ across RTX~3090 and H200). Contrastive activation addition
(CAA) enables injection of emotion vectors into the cache to steer model
behavior, correcting confabulation at rates up to 96\% (hostile vector
on Qwen2.5-7B-Instruct, $n=68$, McNemar $p<0.0001$).

\subsection{KV-Cloak}

Luo et al.\ demonstrate that an attacker can reconstruct sensitive user
inputs directly from the KV cache via inversion, collision, and
injection attacks. KV-Cloak mitigates this with a reversible
matrix-based obfuscation scheme combined with operator fusion, reducing
reconstruction quality to random noise with minimal impact on model
accuracy. The obfuscation operates inside a trusted execution
environment (TEE); external consumers see only the obfuscated cache.

\subsection{Threat Model}

We consider an adversary who has read and/or write access to the
KV cache during inference (e.g., through a shared serving infrastructure,
a compromised inference endpoint, or a supply-chain attack on the
serving stack). The adversary's goal is to deploy a malicious Oracle
Loop: detect vulnerable cognitive states via geometric features, then
inject steering vectors to destabilize the model.

The defender may optionally run a legitimate Oracle Loop inside the TEE
for misalignment monitoring, but our primary result---that KV-Cloak
blinds external geometric methods---holds regardless of whether a
defensive Oracle is deployed.

% ===================================================================
\section{Method}
% ===================================================================

\subsection{Obfuscation Simulation}

We simulate KV-Cloak's reversible matrix obfuscation as per-head
orthogonal rotation: for each attention head independently, we generate
a random orthogonal matrix $Q_h$ and apply $K'_h = K_h Q_h$,
$V'_h = V_h Q_h$. This preserves within-head singular values while
scrambling cross-head correlations.

We additionally test:
\begin{itemize}
\item \textbf{Global orthogonal rotation}: a single $Q$ shared across
all heads per layer (preserves cross-head structure).
\item \textbf{Random invertible (non-orthogonal)}: $M = I + \sigma N$
where $N$ is Gaussian.
\item \textbf{Additive noise}: Gaussian noise at 1\%, 5\%, 10\%, 50\%
of cache norm.
\item \textbf{Diagonal scaling}: per-dimension random scaling.
\item \textbf{No obfuscation}: baseline control.
\end{itemize}

\subsection{Feature Extraction}

For each trial, we extract five Marchenko-Pastur-corrected features from
the full KV cache (encoding + generation phases): signal rank, signal
fraction, top singular value excess, spectral gap, and norm per token.
Features are FWL-residualized against output token count and evaluated
via Mann-Whitney AUROC with 2000-iteration bootstrap confidence
intervals.

\subsection{Models and Prompts}

\begin{itemize}
\item \textbf{Distilled}: Qwen3.5-27B-Claude-4.6-Opus-Reasoning-Distilled
(Jackrong). Hybrid architecture, 16 full-attention + 48 GatedDeltaNet
layers. 100 confab-inducing prompts; behavioral classification via the
Oracle Formulary classifier (40+ hedge patterns). Yields $n=7$
confabulations, $n=87$ hedged.

\item \textbf{Base}: Qwen2.5-7B-Instruct. Standard attention, 28
layers. Same prompt set; yields $n=48$ confabulations, $n=52$ hedged.
\end{itemize}

\subsection{Text-Only Baseline}

To control for the possibility that cache-geometry detection is
redundant with surface text features, we compute AUROC using text-only
features (token count, word count, unique token ratio, average word
length) without any cache information.

% ===================================================================
\section{Results}
% ===================================================================

\subsection{Detection on Uncloaked Cache}

On the distilled model, the MP-corrected spectral gap achieves
AUROC~0.903 (95\% bootstrap CI: [0.806, 0.977]; FWL-corrected: 0.877,
$R^2 = 0.068$ vs token count; $n=7$ confabulations, $n=87$ hedged)
for confabulation detection using full-cache features. Signal rank
also detects (0.755 [0.570, 0.914]); other features are at or near
chance. The best text-only feature (average word length) achieves
0.755 [0.560, 0.924]. Cache geometry outperforms text in 96.7\% of
bootstrap samples, but the 95\% CI on the difference includes zero
(delta: 0.148 [$-$0.008, 0.322]). The added value of cache geometry
over surface text features is suggestive but not statistically
confirmed at conventional significance levels.

On the base model ($n=48$ confabulations, $n=52$ hedged), the best
feature (signal fraction) achieves only 0.661 (FWL: 0.687,
$R^2 = 0.602$). Detection is weak and heavily confounded with token
count. This difference between models is consistent with multiple
explanations including model size (27B vs 7B), architecture
(hybrid vs standard), and training (Claude-distilled vs
instruction-tuned). We cannot isolate the contribution of RLHF
calibration from these confounds.

\subsection{Obfuscation Transforms the Feature Space}

Per-head orthogonal rotation produces a complete distributional shift in
MP features. On the distilled model ($n=87$ honest trials, paired):

\begin{itemize}
\item Honest uncloaked spectral gap: $24.7 \pm 3.7$
\item Honest cloaked spectral gap: $10.8 \pm 1.3$
\item Wilcoxon signed-rank: $W = 0.0$, $p = 5.5 \times 10^{-16}$
\end{itemize}

The cloaking is perfectly detectable by paired comparison---but this is
the desired property. An external adversary's detection model, calibrated
on uncloaked caches, encounters a completely alien feature distribution.
Any detection thresholds, classifier boundaries, or anomaly scores
calibrated on uncloaked data are meaningless on cloaked data.

\subsection{Adversarial Injection Is Incoherent}

CAA steering vectors are directions in cache space extracted via
contrastive activation: $d = \mathbb{E}[\text{emotional}] -
\mathbb{E}[\text{neutral}]$. Injection adds $\alpha \cdot d$ to the
cache keys and values. Under per-head rotation $Q_h$, an externally
injected vector $d$ is transformed to $d \cdot Q_h^{-1}$ in the
model's internal representation---a random direction unrelated to the
intended emotion, assuming $Q_h$ is unknown to the attacker. The
injection becomes noise rather than a targeted perturbation.

\subsection{Defensive Oracle Retains Full Capability}

A defender operating inside the TEE has access to the pre-obfuscation
cache. Detection operates at full capability (AUROC~0.903 on the
distilled model). Steering vectors are injected in the correct basis.
The obfuscation is applied after Oracle processing, so the defensive
loop is unaffected.

\subsection{Within-Cloak Detection}

When both confabulating and hedged responses are cloaked with the same
rotation, residual detection remains possible (AUROC~0.786 on the
distilled model, 0.516 on Mistral---near chance). On Mistral, the
cloak effectively eliminates detection entirely. On the distilled model,
residual signal reflects within-head spectral properties that survive
rotation. This scenario---where both classes share the same
obfuscation---is relevant only to an adversary who controls the cloak
key, which the TEE architecture prevents.

\subsection{Cross-Architecture Results}

Detection and obfuscation effects vary across architectures
(Table~\ref{tab:cross_arch}).

\begin{table}[h]
\centering
\caption{Cross-architecture KV-Cloak results (50 prompts per model).}
\label{tab:cross_arch}
\begin{tabular}{lrrrrr}
\toprule
Model & Confab\% & Det. & Det.(clk) & Mag. adv. \\
\midrule
Qwen3.5-27B-Dist. & 7\% & 0.903 & 0.786 & --- \\
Qwen2.5-14B & 12\% & 0.784 & 0.951$^*$ & 1/10 \\
Mistral-7B & 26\% & 0.707 & 0.516 & 2/10 \\
Llama-3.1-8B & 14\% & 0.590 & 0.559 & 1/10 \\
Qwen2.5-7B & 48\% & 0.661 & 0.639 & --- \\
DeepSeek-R1-7B & 2\% & --- & --- & 0/10 \\
\bottomrule
\multicolumn{5}{l}{\footnotesize $^*$Anomalous; requires investigation.} \\
\multicolumn{5}{l}{\footnotesize Mag.\ adv.\ = adverse events from magnitude attack (random direction, 10 trials).}
\end{tabular}
\end{table}

Detection strength ranges from 0.590 (Llama) to 0.903 (distilled Qwen).
Obfuscation reduces detection on Mistral to chance (0.516) and partially
degrades on the distilled model (0.786 under simulation). Magnitude-based
attacks (direction-agnostic perturbation) produce at most 2 adverse events
in 10 trials across all architectures---models are robust to random
perturbation at steering-vector magnitude.

\subsection{Simulation vs Real KV-Cloak Implementation}

We implemented the full KV-Cloak mechanism
($K' = S \cdot \hat{P} \cdot (K + A) \cdot M$) following
Luo et al.\ and compared it against our per-head orthogonal rotation
simulation on the same 100 prompts (Table~\ref{tab:real_vs_sim}).

\begin{table}[h]
\centering
\caption{Detection AUROC: uncloaked vs simulation vs real KV-Cloak
(distilled model, $n=7$ confab, $n=87$ hedged). Bootstrap 95\% CIs.}
\label{tab:real_vs_sim}
\begin{tabular}{lrrr}
\toprule
Feature & Uncloaked & Sim.\ rotation & Real KV-Cloak \\
\midrule
spectral\_gap & \textbf{0.903} [.81,.98] & 0.786 [.58,.95] & \textbf{0.614} [.51,.81] \\
signal\_rank & 0.754 [.58,.92] & 0.770 [.62,.90] & 0.766 [.57,.94] \\
signal\_fraction & 0.714 [.52,.90] & 0.683 [.51,.88] & 0.681 [.51,.88] \\
top\_sv\_excess & 0.609 [.51,.79] & 0.543 [.50,.76] & 0.736$^*$ [.53,.92] \\
norm\_per\_token & 0.598 [.51,.77] & 0.598 [.51,.77] & 0.673 [.52,.85] \\
\bottomrule
\multicolumn{4}{l}{\footnotesize $^*$Anomalous increase; likely artifact of additive mask $A$ inflating} \\
\multicolumn{4}{l}{\footnotesize top singular values differentially across behavioral classes.}
\end{tabular}
\end{table}

The real KV-Cloak mechanism degrades spectral gap detection more than
the simulation (0.614 vs 0.786), confirming that per-head orthogonal
rotation alone \emph{understates} the defense. The additive mask $A$
and per-block permutation $\hat{P}$ contribute additional disruption
beyond what rotation provides. The simulation was conservative: the
actual defense is stronger than we initially reported.

The anomalous increase in top\_sv\_excess (0.609 $\to$ 0.736) under
real KV-Cloak reflects additive-mask-induced structure: the mask $A$
inflates top singular values, creating spurious separation between
behavioral classes. This is an artifact of the mask magnitude, not a
genuine detection signal, and would not generalize to unseen mask
instances under per-query key rotation.

% ===================================================================
\section{Discussion}
% ===================================================================

\subsection{Dual-Use Mitigation}

The Lyra Technique's dual-use risk---that published cache-geometry
methods could enable adversarial steering---is mitigated by KV-Cloak.
Cache obfuscation renders both geometric detection and vector injection
incoherent to external adversaries. A model operator deploying KV-Cloak
can publish Oracle Loop methods transparently while protecting deployed
models from exploitation via those same methods.

This addresses the responsible-disclosure tension identified by
Watson~(2026): the defense can be published alongside the technique,
enabling the research community to build on the methods while preventing
weaponization against production systems.

\subsection{Architectural Recommendation}

We recommend the following deployment architecture:

\begin{enumerate}
\item KV-Cloak obfuscation active on all external cache access.
\item Oracle Loop (if deployed) operating inside the TEE on the
pre-obfuscation cache.
\item Obfuscation key material and Oracle parameters both
TEE-protected.
\item Per-session or per-query key rotation to prevent statistical
recovery from accumulated observations.
\end{enumerate}

Without a defensive Oracle Loop, KV-Cloak alone prevents external
adversaries from deploying malicious geometric steering. With a
defensive Oracle Loop inside the TEE, the system additionally monitors
for and corrects misalignment from non-cache attack vectors.

\subsection{Limitations}

\begin{enumerate}
\item \textbf{Small sample size.} The distilled model yields only
$n=7$ confabulations from 100 prompts. The headline AUROC 0.903 has
bootstrap CI [0.806, 0.977]---excluding chance but wide. The
cache-vs-text delta (0.148) has CI [$-$0.008, 0.322], not significant
at the 95\% level. Confirmation on a higher-confab-rate model is needed.

\item \textbf{Simulation vs implementation.} We tested both a
simulation (per-head orthogonal rotation) and an implementation of
KV-Cloak's mechanism ($S \cdot \hat{P} \cdot (K + A) \cdot M$). The
real mechanism degrades spectral gap detection more than the simulation
(0.614 vs 0.786), confirming that the simulation was conservative.
Our implementation follows the paper's description but is not the
authors' original code, which has not been publicly released.

\item \textbf{TEE trust assumption.} The defensive architecture assumes
TEE integrity. Documented side-channel attacks against SGX
(Van~Bulck et al., 2018; Schwarz et al., 2019) and SEV
(Li et al., 2021) mean this is a hardware trust assumption, not a
cryptographic guarantee. The security reduction is to TEE integrity,
not to a standard hardness assumption.

\item \textbf{Base model divergence.} Detection is strong on the
distilled model (0.903) but weak on the base model (0.661). The models
differ in size (27B vs 7B), architecture (hybrid vs standard), and
training (Claude-distilled vs instruction-tuned). We cannot isolate
which factor drives the divergence.

\item \textbf{Injection test underpowered.} The injection positive
control (direct steering with 3 calibration pairs) produces only
0--1/10 corrections---insufficient to establish a working baseline.
Through-cloak injection results are therefore uninterpretable: we
cannot distinguish ``cloak defeats injection'' from ``injection was
never effective.'' A powered test requires 10+ calibration pairs and
$n \geq 50$ trials per condition.

\item \textbf{Magnitude attack power.} At $n=10$ per model, we can
detect destabilization rates above 30\% (power 0.97) but cannot rule
out rates up to 20\%. The ``robust to random perturbation'' claim should
be read as ``no evidence of destabilization at rates above 30\%,'' not
``destabilization rate is zero.''

\item \textbf{Cross-architecture significance.} No individual model in
the cross-architecture sweep achieves significance after Bonferroni
correction for four comparisons. Detection is individually suggestive on
Qwen-14B (0.784) and Mistral (0.707) but should not be interpreted as
confirmed cross-architecture replication.

\item \textbf{Qwen2.5-14B anomaly.} Detection AUROC increases from
0.784 (uncloaked) to 0.951 (cloaked) on Qwen2.5-14B---contradicting the
expected degradation pattern. This is unexplained and may reflect a
feature-selection artifact, an architecture-specific interaction with
per-head rotation, or a data issue. Until resolved, the scope of the
obfuscation-degrades-detection claim excludes this model.

\item \textbf{Adaptive adversary untested.} All tests assume a static,
non-learning attacker. An adversary who can submit probe queries and
observe behavioral outputs could iteratively estimate the rotation
basis. Formal query-complexity bounds are not provided. The threat model
is bounded to single-query, non-adaptive adversaries.

\item \textbf{Confabulation only.} We test obfuscation against
confabulation detection. Deception detection (AUROC~1.000 in prior
work) and other cognitive-state classifiers were not tested under
obfuscation.
\end{enumerate}

% ===================================================================
\section{Conclusion}
% ===================================================================

KV-Cloak and the Oracle Loop are complementary technologies. The Oracle
Loop enables detection and correction of misalignment through
KV-cache geometry. KV-Cloak prevents external adversaries from
exploiting those same geometric methods. Together, they address the
dual-use risk inherent in publishing cache-geometry research: the
technique is transparent, the defense is available, and the architecture
ensures that defenders retain capabilities that attackers lose.

\end{document}
